\section{Results}
\label{sec:results}

All numbers are the geometric mean of MASE over the $97$ GIFT-Eval rows, lower
is better. Significance is a two-sided sign test on per-row wins against the
reference. $\star$ marks the two references: \texttt{R0} for v1 arms,
\texttt{D1} for v2 arms.

\begin{table}[htbp]
  \centering\small
  \caption{All trained arms, sorted by MASE. The published \texttt{nano}
  checkpoint scores $1.0742$; our reproduction is \ReproGap\,\% behind it, so
  every comparison below is internal.}
  \label{tab:arms}
  \begin{tabular}{lcclrrrr}
    \toprule
    arm & $k$ & $\beta$ & corpus & MASE & vs \texttt{R0} (\%) & wins/97 & $p$ \\
    \midrule
    + window augmentation & 1 & 1.0 & v2 & 1.1482 & +4.17 & 66 & $5\times10^{-4}$ \\
+ allow\_neg\_eigval & 1 & 1.0 & v2 & 1.1623 & +2.99 & 70 & $2\times10^{-5}$ \\
v2 pipeline & 1 & 1.0 & v2 & 1.1685 & +2.48 & 68 & $9\times10^{-5}$ \\
looped k=4 & 4 & 0.25 & v1 & 1.1729 & +2.11 & 66 & $5\times10^{-4}$ \\
v2 pipeline + looped k=4 & 4 & 0.25 & v2 & 1.1794 & +1.57 & 68 & $9\times10^{-5}$ \\
looped k=8 & 8 & 0.125 & v1 & 1.1864 & +0.99 & 57 & 0.100 \\
looped k=4 & 4 & 0.25 & v1 & 1.1882 & +0.83 & 58 & 0.067 \\
corpus v2 only\,$\star$ & 1 & 1.0 & v2 & 1.1887 & +0.79 & 61 & -- \\
baseline\,$\star$ & 1 & 1.0 & v1 & 1.1982 & +0.00 & -- & -- \\
looped k=2 & 2 & 0.5 & v1 & 1.1996 & -0.12 & 51 & 0.690 \\
looped k=4 beta=1/sqrt k & 4 & 0.5 & v1 & 1.2029 & -0.39 & 56 & 0.150 \\
\bottomrule

  \end{tabular}
\end{table}

\subsection{Looping clears the noise floor by very little}

The two $k=4$ seeds average \LoopMean, which is \LoopGain\,\% better than the
$k=1$ baseline. The two seeds differ from each other by
\NoiseFloorKFour\,\%. The effect is therefore of the same order as the
measurement error: it is positive in both seeds, which is worth something, but
a single-seed claim of ``$+2.1\,\%$ at $k=4$'' would have been an artefact of
reporting the luckier seed.

The sweep over $k$ does not help the case. $k=2$ is indistinguishable from the
baseline ($-0.12\,\%$, $p=0.69$) and $k=8$ gives $+0.99\,\%$ at $p=0.10$ --- so
there is no monotone trend in $k$ to extrapolate, which is what one would
expect if additional iterations were doing additional work. The flow-time
control behaves as predicted: $\beta=1/\sqrt{k}$, which overshoots $T=1$, is
the worst arm in the table.

\subsection{The cost is not small}

\begin{table}[htbp]
  \centering\small
  \caption{Measured training cost of the identical recipe at $k=1$ and $k=4$
  (same corpus, same steps, same hardware). Inference cost scales the same
  way, and is paid on every forecast.}
  \label{tab:cost}
  \begin{tabular}{lcclrr}
    \toprule
    run & $k$ & $\beta$ & steps & wall-clock & Mpts/s \\
    \midrule
    V2A & 1 & 1.0 & 150\,000 & 3\,h\,55 & 11.21 \\
V2L & 4 & 0.25 & 150\,000 & 15\,h\,45 & 2.78 \\
\bottomrule

  \end{tabular}
\end{table}

The wall-clock ratio is \CostRatio$\times$ and the throughput ratio
\ThroughputRatio$\times$, both within $1\,\%$ of the ideal $k$. Looping is
exactly as expensive as stacking $k$ blocks and saves only parameters, which
are not the binding constraint at this scale.

\subsection{Looping does not survive a better pipeline}

\texttt{V2C} and \texttt{V2M} are the cleanest single-variable pair in the
project: same corpus, same sampler, same schedule, same annealing step, $k=1$
versus $k=4$. The looped arm is \VtwoLoopPenalty\,\% \emph{worse}. The
per-horizon split (Table~\ref{tab:perterm}) shows where it goes: $+0.57\,\%$
short, $-4.17\,\%$ medium, $-1.70\,\%$ long. Whatever looping contributed on
the v1 pipeline, a sampler that shows every window exactly once contributes
more, and the two do not add.

\begin{table}[htbp]
  \centering\small
  \caption{Per-horizon decomposition, $55$ short / $21$ medium / $21$ long
  rows. Read within a column only: the three subsets are different data.}
  \label{tab:perterm}
  \begin{tabular}{lrrrr}
    \toprule
    arm & all 97 & short (55) & medium (21) & long (21) \\
    \midrule
    + window augmentation & 1.1482 & 1.1635 & 1.1049 & 1.1524 \\
+ allow\_neg\_eigval & 1.1623 & 1.1886 & 1.1148 & 1.1428 \\
v2 pipeline & 1.1685 & 1.1932 & 1.1165 & 1.1577 \\
looped k=4 & 1.1729 & 1.1975 & 1.1231 & 1.1601 \\
v2 pipeline + looped k=4 & 1.1794 & 1.1864 & 1.1631 & 1.1774 \\
looped k=4 & 1.1882 & 1.2194 & 1.1292 & 1.1682 \\
corpus v2 only & 1.1887 & 1.2137 & 1.1344 & 1.1794 \\
baseline & 1.1982 & 1.2271 & 1.1450 & 1.1778 \\
looped k=4 beta=1/sqrt k & 1.2029 & 1.2298 & 1.1507 & 1.1865 \\
\bottomrule

  \end{tabular}
\end{table}

\subsection{Each cheap lever beats looping}

\begin{table}[htbp]
  \centering\small
  \caption{The three non-architectural levers, each against \texttt{D1}
  (\DOneMase) with one variable changed. Positive is better.}
  \label{tab:levers}
  \begin{tabular}{llrrrr}
    \toprule
    lever & flag & all (\%) & short (\%) & medium (\%) & long (\%) \\
    \midrule
    window augmentation & \texttt{--augment} & +3.41 & +4.14 & +2.60 & +2.29 \\
allow\_neg\_eigval & \texttt{--allow-neg-eigval} & +2.22 & +2.07 & +1.73 & +3.10 \\
v2 pipeline (exact-once epochs) & \texttt{--sampler enumerated} & +1.70 & +1.69 & +1.58 & +1.84 \\
\bottomrule

  \end{tabular}
\end{table}

Augmentation gives $+3.41\,\%$, the negative-eigenvalue gate $+2.22\,\%$, the
exact-once sampler $+1.70\,\%$ --- against looping's \LoopGain\,\% at
\CostRatio$\times$ the compute. The first and third are essentially free; the
second costs nothing at inference and, per \citet{ckda2026}, $3$--$4\,\%$ of
kernel throughput.

\subsection{Only the eigenvalue lever behaves like a representational change}

The per-horizon columns of Table~\ref{tab:levers} separate the three levers by
\emph{mechanism}, which is the result we consider most interesting.
Augmentation helps most at short horizon and decays with length
($+4.14 / +2.60 / +2.29$), the signature of regularisation: it improves each
individual patch, and the gain dilutes along the autoregressive rollout. The
sampler is flat ($+1.69 / +1.58 / +1.84$) --- it simply trains on more of the
data. Only \texttt{allow\_neg\_eigval} \emph{grows} with horizon
($+2.07 / +1.73 / +3.10$), and it is the only lever for which that pattern was
predicted in advance, from the argument in \S\ref{sec:background} that
reflections are what let a linear recurrence carry alternating structure across
many steps.

This is a single arm and one should not over-read it. But it is the one place
in the project where a pre-registered mechanistic prediction was confirmed, and
it points at the spectrum of the transition rather than the depth of the stack.

\subsection{Composition (in progress)}

Two seeds of the composed arm \texttt{AUGNEG} (augmentation $+$
\texttt{allow\_neg\_eigval}, v2 corpus, $k=1$) have finished training; the
evaluations are pending. If the two levers are independent, the multiplicative
prediction is $\AugNegPrediction$; adding the v2 sampler would predict
$\TripleLeverPrediction$. The data streams were verified identical to
\texttt{AUG}'s by \texttt{stream\_sha256} and the model flag verified set, so
the arm tests composition and not an accidental third difference.

\todo{as soon as a GPU frees}{%
  \texttt{AUGNEG} / \texttt{AUGNEGS1} benchmark numbers, and the full
  $97$-row CRPS of the quantile head (\texttt{R1Q}), whose smoke test on
  \texttt{m4\_weekly} gave $-11.1\,\%$ CRPS at unchanged MASE.}
